%!TEX root = ../Pflichtenheft.tex

\chapter{Definition of Objectives}\label{ch:definition-of-objectives}

The initial public release of \textit{ChatGPT} by \textit{OpenAI} in November 2022 was the start of publicly available and free generative \hyperref[abr:AI]{AI} products.
Since then, a variety of \textit{Large Language Models (\hyperref[abr:LLM]{LLM}s)} have been released and updated by a number of different organizations.
Due to their extensive availability, these models have been rapidly adopted in a variety of applications across many facets of human life.
As \hyperref[abr:LLM]{LLM}s have become more powerful, one common use case of them has become the automatic generation of source code based on natural language input. \newline
Generating code samples of small to medium-sized code snippets for prominent languages like \textit{Java} or \textit{C} is an effective strategy that is widely employed.
However, \textit{\hyperref[abr:LLM]{LLM}s} continue to encounter challenges in handling \textit{Domain Specific Languages (\hyperref[abr:DSL]{DSL}s)}. \textit{\hyperref[abr:DSL]{DSL}s} are programming languages, configuration formats, or embeddings into more common programming languages that are only used in very specific scenarios or industries.
Common examples of niche \textit{\hyperref[abr:DSL]{DSL}s} are \textit{nginx configurations}, \textit{R}, or \textit{Gradle scripts}. \newline
This is particularly evident for \textit{\hyperref[abr:DSL]{DSL}s} that are used for hardware description, \hyperref[abr:PLC]{PLC} control and low-level machine programming.
Due to being specific for hardware, proprietary, or limited to products of a company, these languages are only sparsely included in the training data for \textit{\hyperref[abr:LLM]{LLM}s}. \newline
Solving this issue is the objective of the \textbf{CodeX} system. \textbf{CodeX} aims to allow users to generate \textit{\hyperref[abr:DSL]{DSL}} source code from natural language, specifically the \textit{\hyperref[abr:SPICE]{SPICE} netlist} syntax that is used to describe electronic circuits.
It will also assist the user to interpret the generated circuit description by providing visualization.
The backing \textit{\hyperref[abr:LLM]{LLM}} will be supported by a database of examples of that \textit{\hyperref[abr:DSL]{DSL}}.


\section{Must Criteria}\label{sec:must-criteria}
\textbf{Must Criteria} describe aspects of the system that have to be present in order for the program to be usable on a basic level.

\must{1}{The system must provide a web interface that allows natural language to be converted into a snippet of code. The target \textit{\hyperref[abr:DSL]{DSL}} can be selected out of options and the generated code can be exported as a file.}
\must{2}{The system must provide \textit{\hyperref[abr:SPICE]{SPICE}} as one option of its selectable \textit{\hyperref[abr:DSL]{DSL}s}.}
\must{3}{\textit{\hyperref[abr:SPICE]{SPICE}} conversations must have a visualization of the respective circuit.}
\must{4}{Every conversation must have a name.}
\must{5}{The system must support the requests to the \textit{\hyperref[abr:LLM]{LLM}} using a \textbf{\hyperref[abr:RAG]{RAG}} approach. The data that is retrieved must be stored on a secure database that respects \textbf{\hyperref[abr:IP]{IP} property}.}
\must{6}{The system must integrate a validation back-end to validate the generated circuit for \textit{\hyperref[abr:SPICE]{SPICE}} conversations.}
\must{7}{The system must allow adding/removing separate \textit{\hyperref[abr:LLM]{LLM}} models together with their \textit{OpenAI API} URL. When a request is made, the user should be able to select the model with which the request will be performed.}

\section{Should Criteria}\label{sec:should-criteria}
\textbf{Should criteria} describe aspects of the system that are not essential to a basic usage of the system but should be fulfilled to meet the project's initial objectives.

\should{1}{The circuit visualization for \textit{\hyperref[abr:SPICE]{SPICE}} should be interactive.}
\should{2}{The interface of the system should allow one to cancel a pending request.}
\should{3}{Users should be able to sign up, log in and log out of accounts using a username (or email) and password.}
\should{4}{A strict password complexity policy should be implemented.}
\should{5}{The system should allow users to securely update their profile data.}
\should{6}{The system should provide a secure mechanism for users to recover their accounts in case of lost credentials.}
\should{7}{Conversations are saved with respect to the account that created them and are only visible to that account.}
\should{8}{If the \textit{\hyperref[abr:LLM]{LLM}} fails to generate working code after multiple retries, the system should stop further attempts and display an error message indicating the failure to generate the requested code.}
\should{9}{The system should provide a list of all previous conversations and allow to view them.}
\should{10}{Every conversation should be able to be marked as deleted and won't be shown in the overview after that.}

\newpage
\section{Could Criteria}\label{sec:could-criteria}
\textbf{Could Criteria} describe aspects of the system that add onto the originally specified targets of the system and are not required.

\could{1}{The system could integrate a simulation back-end to simulate the generated circuit for \textit{\hyperref[abr:SPICE]{SPICE}} conversations.}
\could{2}{The system could accept follow-up instructions on conversations.}
\could{3}{The system could provide more \textit{\hyperref[abr:DSL]{DSL}} options than only \textit{\hyperref[abr:SPICE]{SPICE}}.}
\could{4}{Every account could be associated to one of the roles \textbf{\hyperref[def:admin]{Admin}} or \textbf{\hyperref[def:member]{Member}}. Every role could have a specified upper bound of requests per month than could be set as a server parameter.}
\could{5}{The role request upper bounds could be set from admins via the web interface.}
\could{6}{The web interface could be offered in multiple languages, at least \textit{German} and \textit{English}.}
\could{7}{If the \textit{DSL} allows it, code output could be prettified using syntax highlighting.}
\could{8}{The front-end could allow users to pin a conversation so that it remains at the top of the conversation list.}
\could{9}{During registration, the user's email could be confirmed by sending a verification link or code. Registration would not be completed until the email is confirmed. When the user would change their email, a new email confirmation would automatically be issued.}
\could{10}{The system could support \hyperref[def:oauth]{OAuth} 2.0 authentication to allow users to log in easily using existing accounts from trusted providers, such as Google.}
\could{11}{The software could provide a dark mode as an alternative theme.}
\could{12}{The system could integrate with cloud storage providers to save or share generated code snippets.}
\could{13}{The app could allow users to search conversations by name to quickly locate specific conversations.}
\could{14}{The front-end could indicate to the user that a request is currently processing.}
\could{15}{The software could allow users to download generated code as a file.}
\could{16}{The software could provide admin-only analytics on usage, such as daily active users, which are not visible to regular users.}
\could{17}{Users could upload and display a profile picture associated with their account.}

\newpage

\section{Won't Criteria}\label{sec:wont-critaria}
\textbf{Won't Criteria} describe aspects that are not planned to be included in the system.

\wont{1}{A user manual won't be delivered.}
\wont{2}{The circuit visualization for \textit{\hyperref[abr:SPICE]{SPICE}} conversations won't allow editing the circuit.}
\wont{3}{Users won't be able to attach any more data to a request than their prompt.}
\wont{4}{The application won't support real-time collaboration between multiple users in the same conversation.}
\wont{5}{The application won't provide another client than the web-client.}
\wont{6}{Voice input or voice-based commands won't be included in the software.}
\wont{7}{The application won't use interactions with users to train or improve the \textit{\hyperref[abr:LLM]{LLM}}.}